% !Mode:: "TeX:UTF-8"
% !TEX program = xelatex
% 天津科技大学本科生毕业设计（论文）模板宏包依赖
% 修改整理：刘建征 (2024.4.20)

% =============================================
% 基础版面与语言环境
% =============================================
\usepackage[a4paper,text={165.0true mm,242true mm},top=25.0true mm,left=25.0true mm,head=7.5true mm,headsep=7.5true mm,foot=17true mm]{geometry}
\usepackage{setspace}       % 行距控制
\usepackage{indentfirst}    % 首行缩进宏包

% =============================================
% 数学、物理符号与字体
% =============================================
\usepackage{amsmath}        % AMSLaTeX宏包，用来排出更加漂亮的公式
\usepackage{amssymb}        % 数学符号生成命令
\usepackage{bm}             % 处理数学公式中的黑斜体的宏包
\usepackage[amsmath,thmmarks,hyperref]{ntheorem} % 定理类环境宏包
\usepackage[squaren]{SIunits} % 支持国际标准单位
\usepackage{fontspec}       % 字体配置宏包

% =============================================
% 图表控制
% =============================================
\usepackage{graphicx}       % 支持插图处理
\usepackage[below]{placeins}% 浮动图形控制
\usepackage{subcaption}     % 现代子图宏包，替代老旧的 subfigure 和 ccaption
\usepackage{booktabs}       % 三线表宏包
\usepackage{multirow}       % 允许表格合并多个单元格
\usepackage{longtable}      % 支持跨页的长表格
\usepackage{tabularx}       % 自动设置表格的列宽
\usepackage{array}          % 扩展的表格列格式

% =============================================
% 代码与算法环境
% =============================================
\usepackage{minted}         % Pygments 语法高亮，编译需 -shell-escape
\usepackage[boxed,ruled,lined]{algorithm2e} % 算法排版
\usepackage{algorithmic}

% =============================================
% 参考文献与交叉引用
% =============================================
\usepackage[sort&compress,numbers]{natbib}  % 参考文献格式
\usepackage{xurl}           % 允许长URL自动换行

% =============================================
% 杂项与工具
% =============================================
\usepackage{xcolor}         % 颜色支持
\usepackage{pagecolor}      % 背景颜色控制
\usepackage{tikz}           % 强大的绘图工具
\usepackage{enumitem}       % 改变列表项的格式
\usepackage{calc}           % 支持长度的计算
\usepackage{ulem}           % 下划线支持
\usepackage{adjustbox}      % 盒子与图像缩放调整

% =============================================
% 目录控制
% =============================================
\usepackage{titletoc}       % 控制目录的宏包

% =============================================
% 超链接设置 (建议放在最后加载)
% =============================================
\usepackage[unicode,
            pdfstartview=FitH,
            bookmarksnumbered=true,
            bookmarksopen=true,
            colorlinks=false,
            pdfborder={0 0 1},
            citecolor=blue,
            linkcolor=red,
            anchorcolor=green,
            urlcolor=blue,
            breaklinks=true
            ]{hyperref}